% Иллюстрация к выводу статвеса бозонов:
% N бозонов в Z ячейках, разделённых Z-1 перегородками; заполнение произвольное.
% Компиляция: pdflatex bose-cells-figure.tex
\documentclass[tikz,border=8pt]{standalone}
\usepackage[T2A]{fontenc}
\usepackage[utf8]{inputenc}
\usepackage[russian]{babel}
\usetikzlibrary{arrows.meta,decorations.pathreplacing}

\begin{document}
\begin{tikzpicture}[
    note/.style = {gray!30!black},
    ptr/.style  = {gray!70, thin},
  ]

  % ================= ячейки =================
  % дно
  \draw[semithick] (0,0) -- (10,0);
  % внешние стенки (не считаются перегородками)
  \draw[line width=1.2pt] (0,0) -- (0,1.12);
  \draw[line width=1.2pt] (10,0) -- (10,1.12);
  % Z-1 = 5 перегородок
  \foreach \x in {1.667,3.333,5,6.667,8.333} {
    \draw[semithick] (\x,0) -- (\x,1.0);
  }

  % ================= бозоны: заполнение произвольное =================
  % занятость ячеек: 3, 0, 2, 1, 0, 2  (N = 8)
  \foreach \x in {0.40,0.84,1.28, 3.90,4.45, 5.83, 8.90,9.45} {
    \shade[ball color=gray!45] (\x,0.23) circle (0.21);
    \draw[gray!60, thin] (\x,0.23) circle (0.21);
  }

  % ================= подписи =================
  \node[note] at (1.5,1.9) {$N$ бозонов};
  \draw[ptr] (1.36,1.60) -- (0.86,0.52);
  \node[note] at (5.9,1.9) {$Z-1$ перегородок};
  \draw[ptr] (5.42,1.66) -- (5.02,1.06);

  \draw[decorate, decoration={brace, mirror, amplitude=6pt}]
        (0,-0.28) -- (10,-0.28) node[midway, below=8pt] {$Z$ ячеек};

\end{tikzpicture}
\end{document}
